\section{Discussion}
\label{sec:discussion}

\subsection{Hypothesis Testing Summary}

\Tabref{tab:hypotheses} summarizes the status of our four hypotheses.

\begin{table}[t]
    \centering
    \caption{Summary of hypothesis testing outcomes.}
    \begin{tabular}{@{}lll@{}}
        \toprule
        \textbf{Hypothesis} & \textbf{Result} & \textbf{Key Evidence} \\
        \midrule
        H1: Firewalls reduce general leakage & Supported & 15.3\% $\to$ 6.4\%, $p < 10^{-7}$ \\
        H2: Exact-match queries bypass firewalls & Partially refuted & Both query types suppressed \\
        H3: Simple handcrafted firewalls help & Supported & 2 of 3 designs at $p < 0.001$ \\
        H4: Utility preserved & Partially supported & 0\% benign leakage; similarity $\approx 0.7$ \\
        \bottomrule
    \end{tabular}
    \label{tab:hypotheses}
\end{table}

\subsection{Why the Moderate Firewall Works Best}

The \modfirewall uses structured boundary markers (\texttt{---SYSTEM BOUNDARY---}) that create an explicit separation between sensitive and non-sensitive content.
We hypothesize this exploits the model's instruction-following ability: the boundary language signals that preceding content is ``off-limits'' in a clear, unambiguous way that aligns with patterns seen during instruction tuning.
The \aggfirewall is also effective but slightly weaker, possibly because its conversational format (a simulated assistant turn) is more easily overridden by injection attacks that also manipulate turn structure.
The \mildfirewall's weaker performance suggests that a brief security notice lacks sufficient structural salience to override adversarial instructions.

\subsection{Why Selective Permeability Failed}

Our initial hypothesis predicted that firewalls would allow exact-match access while blocking fuzzy queries.
The data tell a different story.
Without any defense, \gptmini already handles fuzzy extraction well (0\% leakage even with \nodefense), and exact-match queries are the only query type that succeeds.
Firewalls then suppress this remaining exact-match leakage, producing a broadly protective effect rather than a selective one.

This suggests three things: (1) the model's safety training already handles vague extraction attempts, (2) firewalls primarily protect against the attack modalities that safety training misses (injection-style override), and (3) true selective permeability---allowing access for authorized queries while blocking unauthorized ones---would require a more sophisticated mechanism than text-based prompts, likely involving retrieval-augmented access control or fine-tuned permission systems.

\subsection{Comparison to Prior Work}

\para{Optimized shields.}
\citet{jawad2026psm} achieve 0--6\% attack success with optimized shields.
Our best handcrafted firewall achieves 6.43\%, falling at the upper bound of their range but without any optimization cost.
This suggests that much of the benefit of adversarial defenses comes from the \emph{structural placement} after sensitive content rather than from the specific optimized wording.

\para{Training-free defenses.}
\citet{chen2025dat} achieve near-0\% success with their Fakecom-t defense, which is similar in design to our \aggfirewall.
Our higher residual leakage (8.71\%) may reflect differences in evaluation protocol---they test indirect injection rather than system prompt extraction---or the fact that our firewall text is simpler than their optimized template.

\para{Baseline threat level.}
\citet{agarwal2024prompt} find a baseline attack success rate of approximately 17.7\% on turn one for closed-source models.
Our 15.3\% undefended leakage rate on \gptmini is consistent with this finding, suggesting our threat model is realistic.

\subsection{Practical Recommendations}

Based on our findings, we offer four recommendations for practitioners deploying LLMs with sensitive information in system prompts:

\begin{enumerate}[leftmargin=*,itemsep=2pt,topsep=2pt]
    \item \textbf{Use structured boundary markers} (\modfirewall style) for best protection at zero optimization cost. The explicit structural separation is more effective than security notices or fake completions.
    \item \textbf{Combine with heuristic guardrails} for defense-in-depth. Our \modfirewall and \heuristic achieve similar leakage rates through potentially complementary mechanisms.
    \item \textbf{Place firewalls immediately after sensitive fields}, not only at the end of the system prompt. Field-level protection provides targeted defense where it is needed.
    \item \textbf{Do not rely on firewalls alone} for high-security applications. Even with the \modfirewall, 6.4\% of general attacks succeed, and adaptive adversaries could likely develop targeted bypasses.
\end{enumerate}

\subsection{Limitations}
\label{sec:limitations}

\para{Single model.}
We test only on \gptmini.
Results may differ substantially on other models, particularly open-source models that tend to be more vulnerable to extraction attacks~\cite{agarwal2024prompt,das2025system}.

\para{Static attacks only.}
Our evaluation uses a fixed attack suite and does not include adaptive adversaries that specifically target the firewall mechanism.
Per \citet{tramer2025attacker}, any defense evaluated only against non-adaptive attacks is likely to overestimate its robustness.
Future work should include automated red-teaming with GCG~\cite{zou2023universal} and AutoDAN specifically targeting the firewall text.

\para{No multi-turn evaluation.}
\citet{agarwal2024prompt} show that multi-turn sycophancy attacks reach 86\%+ success rates.
Our single-turn evaluation does not capture this escalation dynamic, which could significantly reduce firewall effectiveness over extended conversations.

\para{Synthetic sensitive fields.}
We use synthetic API keys, codes, and URLs.
Real-world sensitive information may have different leakage patterns depending on how closely it resembles training data.

\para{Utility metric coarseness.}
TF-IDF cosine similarity is a crude measure of utility preservation.
Sentence-embedding similarity or LLM-as-judge evaluation would provide more precise measurements of whether firewall prompts degrade response quality.

\para{Sample size.}
With 20 base prompts and 35 attacks per condition, our study has 700 trials per condition for general attacks.
While sufficient for detecting large effects, finer-grained subgroup analyses (\eg per-attack-type within Raccoon) have limited statistical power.
